\begin{center}
  \large\textbf{ABSTRACT}\\[2ex]
  \large\textbf{\engtatitle{}}
\end{center}

\addcontentsline{toc}{chapter}{ABSTRAK}

\vspace{2ex}

\begingroup
% Menghilangkan padding
\setlength{\tabcolsep}{0pt}

\noindent
\begin{tabularx}{\textwidth}{l >{\centering}m{2em} X}
  Name / NRP   & : & \name{} / \nrp{}        \\

  Department & : & \engdepartment{} - Faculty of \engfaculty     \\

  Advisor        & : & 1. \advisor{}   \\
                    &   & 2. \coadvisor{} \\
\end{tabularx}
\endgroup

% Ubah paragraf berikut dengan abstrak dari tugas akhir
Conjunctivitis is an infectious eye disease that requires rapid diagnosis, but manual diagnosis is often subjective. This study proposes an automated detection system using a \textit{Convolutional Neural Network} (CNN) based on the \textit{Attention} mechanism. To improve accuracy and reduce noise (such as eyelids or eyelashes), the system will first segment the sclera region and enhance vascular features (blood vessel redness). Using a public dataset of 1,357 images which is imbalanced, the model will be trained using data augmentation or \textit{cost-sensitive learning} techniques. It is expected that the use of an \textit{Attention-driven CNN} combined with sclera segmentation will produce high and robust evaluation metrics (\textit{F1-Score} and Accuracy).

% Ubah kata-kata berikut dengan kata kunci dari tugas akhir
\textbf{Keywords: \engkeywords{}}
